\section{Conclusion}
In conclusion, both models have their strengths and weaknesses. Sverdrup theory is a good intuitive way to visualize the flow in the mid ocean, while Munk theory is good if one also wants to resolve the western boundary current. Three important conclusions are:
\begin{itemize}
    \item Both numerical models showed western boundary intensification, explained by Rossby wave propagation. However these Rossby waves must be balanced by some kind of friction to die with time and not cause non-converging behaviour.
    \item Although both numerical models showed return flow, only Munk theory showed a western boundary current, explained by a frictional boundary layer.
    \item A larger domain gave the same interior meridional velocity, which resulted in a stronger western boundary current due to the larger flux of the larger ocean.
\end{itemize}
One can summarize the two models by stating that although the Sverdrup theory is more intuitive from a physical perspective, the artificial Munk theory presented a more realistic solution. However, a good physical-oceanographer would then know when to use one model or the other, in order to explain different aspects of the wind driven ocean gyres.